\section*{Bab \ref{ch:g11:ocean-ridges} --- Punggung Samudra dan Litosfer Samudra}

\begin{solution}{exo:g11:ocean-ridges:1}
\omterm{def:g11:ocean-ridges:peridotite}{Peridotit} adalah batuan gelap berbutir kasar yang sebagian besar terdiri atas
olivin, dengan sedikit piroksen; sebagian besar mantel terdiri atas batuan ini.
\omterm{def:g11:ocean-ridges:gabbro}{Gabro} adalah batuan beku gelap berbutir kasar yang terdiri atas feldspar
plagioklas dan piroksen, batuan yang paling melimpah di \omterm{def:g11:ocean-ridges:oceanic-lithosphere}{kerak samudra} bagian
dalam. \omterm{def:g11:ocean-ridges:gabbro}{Gabro} termasuk kerak, \omterm{def:g11:ocean-ridges:peridotite}{peridotit} termasuk mantel. \omterm{def:g11:ocean-ridges:gabbro}{Gabro} mengkristal dengan
lambat, jauh di bawah dasar laut, tempat magma mendingin dengan lambat:
kristal-kristalnya mempunyai waktu untuk tumbuh. Lava basalt mendingin dengan
cepat di dasar laut, dan kristal-kristalnya tetap kecil.
\end{solution}

\begin{solution}{exo:g11:ocean-ridges:2}
Dari dasar laut ke bawah: sedimen (tidak ada di sumbu), lava basalt (lava bantal
dan aliran lembaran), retas berlapis, \omterm{def:g11:ocean-ridges:gabbro}{gabro}, lalu, di bawah \omterm{def:g11:seismic-structure:moho}{Moho}, \omterm{def:g11:ocean-ridges:peridotite}{peridotit}
mantel. Lubang itu menembus $1507/7100 \approx 0.21$, sekitar seperlima
kerak: \omterm{def:g11:ocean-ridges:gabbro}{gabro} bagian dalam dan \omterm{def:g11:seismic-structure:moho}{Moho} belum pernah dicapai oleh bor.
\end{solution}

\begin{solution}{exo:g11:ocean-ridges:3}
$d(0) = \qty{2500}{m}$; $d(9) = 2500 + 350 \times 3 = \qty{3550}{m}$;
$d(25) = 2500 + 350 \times 5 = \qty{4250}{m}$.
\end{solution}

\begin{solution}{exo:g11:ocean-ridges:4}
\omterm{def:g11:ocean-ridges:normal-fault}{Sesar normal} adalah sesar miring yang di sepanjangnya blok di atas bidang sesar
telah bergerak turun relatif terhadap blok di bawahnya. Di punggung yang lambat,
lempeng-lempeng menarik \omterm{def:g11:lithosphere-asthenosphere:lithosphere}{litosfer} yang dingin dan getas sehingga terpisah:
\omterm{def:g11:lithosphere-asthenosphere:lithosphere}{litosfer} itu teregang, dan patah. Blok-blok di atas sesar meluncur turun, ke arah
sumbu, dan lembah sumbu adalah blok yang telah turun di antara sesar-sesar itu.
\end{solution}

\begin{solution}{exo:g11:ocean-ridges:5}
Air itu tidak mendidih karena tekanan: lebih dari dua kilometer air menekannya,
dan di bawah tekanan itu air tetap cair jauh di atas \qty{100}{\celsius}.
``Asap'' hitam itu adalah awan butir-butir halus besi sulfida dan
mineral-mineral lain, yang mengendap ketika air ventilasi yang panas, yang kaya
akan logam terlarut, bertemu dengan air laut yang dingin.
\end{solution}

\begin{solution}{exo:g11:ocean-ridges:6}
(1) Pada kedalaman sekitar \qty{54}{km}, pada suhu sekitar \qty{1324}{\celsius}.
(2) Pada kedalaman sekitar \qty{85}{km}, pada suhu sekitar
$1400 + 0.45 \times 85 \approx
\qty{1438}{\celsius}$. (3) Mantel yang lebih panas mulai meleleh lebih dalam,
sehingga kolom batuan yang naik yang meleleh sebelum mencapai kerak lebih
panjang: mantel itu membuat lebih banyak basalt, dan kerak yang lebih tebal.
\end{solution}

\begin{solution}{exo:g11:ocean-ridges:7}
(1) $1489 - 1345 = \qty{144}{\celsius}$ di bawah suhu lelehnya.
(2) Batuan itu naik $100 - 54 = \qty{46}{km}$; selisihnya menyempit rata-rata
$144/46 \approx \qty{3.1}{\celsius}$ per kilometer. (3) Tidak: selama naik batuan
itu bahkan sedikit mendingin, sekitar $0.45 \times 46 \approx
\qty{21}{\celsius}$. Suhu lelehnyalah yang turun, jauh lebih cepat, ketika
tekanan berkurang.
\end{solution}

\begin{solution}{exo:g11:ocean-ridges:8}
Laju setengahnya \qty{7.3}{cm} dan \qty{1.65}{cm} setahun, yaitu
\qty{73}{km} dan \qty{16.5}{km} per sejuta tahun. Untuk membuat dasar sepanjang
\qty{1000}{km} di satu sisi: $1000/73 \approx 13.7$ juta tahun untuk Tinggian
Pasifik Timur, $1000/16.5 \approx 61$ juta tahun untuk Punggung Tengah Atlantik.
\end{solution}

\begin{solution}{exo:g11:ocean-ridges:9}
(1) \qty{5}{cm} sama dengan \qty{5e-5}{km}: $65000 \times 5 \times 10^{-5} =
\qty{3.25}{km^2}$ dasar baru setahun. (2) $3.25 \times 7.1 \approx
\qty{23}{km^3}$ kerak setahun. (3) $3.6 \times 10^8 / 3.25 \approx 1.1
\times 10^8$ tahun, sekitar 110 juta tahun: dengan laju ini punggung-punggung
samudra dapat membuat seluruh dasar samudra dalam sekitar seratus juta tahun,
waktu yang singkat dibandingkan dengan umur Bumi. Itulah sebabnya tidak ada dasar
samudra yang sangat tua.
\end{solution}

\begin{solution}{exo:g11:ocean-ridges:10}
(1) Bidang sesar membentuk sudut \ang{60} dengan bidang datar, dan penurunannya
adalah sisi tegak segitiga siku-siku yang sisi miringnya adalah pergeseran itu:
$1.5/\sin\ang{60} \approx \qty{1.73}{km}$. (2) Setiap sesar menjauhkan kedua sisi
secara mendatar sejauh $1.5/\tan\ang{60} \approx \qty{0.87}{km}$; kedua sesar
bersama-sama, sekitar \qty{1.73}{km}. (3) Tidak. \qty{25}{km} itu adalah jarak
antara sesar-sesar yang membatasi blok yang turun, bukan besarnya peregangan,
yang hanya sekitar \qty{1.7}{km}; dan punggung itu terus membuka oleh retas dan
lava baru di sumbunya selain oleh pensesaran.
\end{solution}

\begin{solution}{exo:g11:ocean-ridges:11}
$d(25) = 6400 - 3200\,e^{-25/62.8} \approx 6400 - 3200 \times 0.672 \approx
\qty{4251}{m}$; $d(100) = 6400 - 3200\,e^{-100/62.8} \approx 6400 - 3200
\times 0.203 \approx \qty{5749}{m}$. Antara 25 dan 100 juta tahun dasar itu
tenggelam sekitar \qty{1500}{m}. Hukum akar kuadrat memberikan
$2500 + 350 \times 5 = \qty{4250}{m}$ pada 25 juta tahun: kedua hukum itu
sesuai, di tempat keduanya berlaku.
\end{solution}

\begin{solution}{exo:g11:ocean-ridges:12}
(1) $\alpha\,\Delta T\,L = 3.2 \times 10^{-5} \times 700 \times 100 =
\qty{2.24}{km}$. (2) $s = 3380 \times 2.24/(3380 - 1025) = 7571/2355 \approx
\qty{3.2}{km}$. (3) Hukum eksponensial memungkinkan penenggelaman sebesar
$6400 - 2500 =
\qty{3900}{m}$ seluruhnya: pendinginan dan penyusutan saja menjelaskan sekitar
empat perlimanya, dengan nilai-nilai sederhana untuk penurunan suhu dan tebalnya.
\end{solution}

\begin{solution}{exo:g11:ocean-ridges:13}
(1) $54 - 7 = \qty{47}{km}$ dibandingkan dengan $85 - 7 = \qty{78}{km}$: kolom
yang lebih panas $78/47 \approx 1.66$ kali lebih panjang. (2) $7.1 \times 1.66 \approx
\qty{11.8}{km}$. (3) \qty{10.3}{km} yang terukur lebih tebal daripada normal,
seperti yang diramalkan, tetapi kurang dari taksiran itu: kilometer-kilometer
terdalam kolom itu, tepat di atas awal pelelehan, memberikan lelehan yang lebih
sedikit daripada bagian yang lebih dangkal, sehingga anggapan lelehan yang sama
per kilometer memberikan taksiran yang terlalu besar.
\end{solution}

\begin{solution}{exo:g11:ocean-ridges:14}
(1) Hukum itu memberikan \qty{3178}{m}, \qty{4160}{m}, \qty{4782}{m}, dan
\qty{5267}{m}; median yang terukur lebih dangkal sebesar \qty{64}{m},
\qty{142}{m}, \qty{317}{m}, dan \qty{619}{m}. (2) Selisihnya bertambah seiring
umur. (3) Hukum itu memberikan kedalaman dasar itu sendiri, yang dikoreksi untuk
sedimennya; sel-sel yang terukur membawa sedimennya, yang lebih tebal di dasar
yang lebih tua, dan mencakup gunung bawah laut dan dataran tinggi vulkanik, yang
lebih dangkal daripada dasar di sekitarnya.
\end{solution}

\begin{solution}{exo:g11:ocean-ridges:15}
(1) $p = 1025 \times 9.8 \times 2500 \approx \qty{25.1}{MPa}$, ditambah
\qty{0.1}{MPa}: sekitar \qty{25.2}{MPa}. (2) Nilai ini di atas tekanan kritis,
\qty{22.064}{MPa}: air ventilasi itu tidak dapat mendidih, betapapun panasnya.
(3) Tekanan mencapai \qty{22.064}{MPa} pada kedalaman
$(22.064 - 0.1) \times 10^6/(1025 \times 9.8) \approx \qty{2190}{m}$: setiap
ventilasi yang lebih dalam daripada sekitar \qty{2200}{m} aman dari pendidihan.
\end{solution}

\begin{solution}{pb:g11:ocean-ridges:1}
\textbf{1.} $14.6/2 = \qty{7.3}{cm}$ setahun, yaitu \qty{73}{km} per sejuta
tahun.

\textbf{2.} Pada kedalaman sumbu punggung, $d(0) = \qty{2500}{m}$.

\textbf{3.} Dari dasar laut ke bawah: lava basalt (lava bantal dan aliran
lembaran), retas berlapis, dan \omterm{def:g11:ocean-ridges:gabbro}{gabro}, seluruhnya kerak setebal sekitar
\qty{7.1}{km}; di bawah \omterm{def:g11:seismic-structure:moho}{Moho}, \omterm{def:g11:ocean-ridges:peridotite}{peridotit} mantel, sampai dasar \omterm{def:g11:lithosphere-asthenosphere:lithosphere}{litosfer}, yang
masih sangat tipis di sumbu. Sedimen akan mengendap di atasnya kemudian.

\textbf{4.} Ketika lempeng-lempeng saling menjauh, \omterm{def:g11:lithosphere-asthenosphere:lithosphere}{astenosfer} naik di bawah
punggung. Suhunya hanya turun sekitar \qty{0.45}{\celsius} per kilometer,
sedangkan suhu lelehnya turun jauh lebih cepat ketika tekanan berkurang. Pada
kedalaman sekitar \qty{54}{km} keduanya bertemu: batuan mulai meleleh, dan terus
meleleh selama naik, tanpa ada kalor yang ditambahkan (\omterm{def:g11:ocean-ridges:partial-melting}{pelelehan dekompresi}).

\textbf{5.} Setiap sejuta tahun punggung itu membuat dasar sepanjang
$2 \times 73 = \qty{146}{km}$ di kedua sisinya; untuk setiap kilometer punggung,
volume keraknya $146 \times 7.1 \times 1 \approx \qty{1040}{km^3}$.

\textbf{6.} $73 \times t$: \qty{73}{km}, \qty{292}{km}, \qty{1168}{km}, dan
\qty{4672}{km}.

\textbf{7.} $2000/73 \approx 27$ juta tahun.

\textbf{8.} Lempeng-lempeng mengubah kecepatan dan arahnya dalam puluhan juta
tahun, seperti yang ditunjukkan oleh tikungan rantai Hawaii--Emperor
(\cref{ch:g11:plate-motions}); laju masa kini tidak harus berlaku selama 64
juta tahun.

\textbf{9.} Jauh dari sumbu kerak tidak lagi berada di atas lensa magma: kerak
itu mendingin, dan di bawahnya tidak ada lagi cukup kalor untuk menggerakkan
sirkulasi dan memanaskan air.

\textbf{10.} Di barat Tinggian Pasifik Timur potongan itu menempuh perjalanan di
atas lempeng Pasifik. Di sebelah barat, lempeng Pasifik berakhir di sepanjang
palung-palung dalam, tempat lempeng itu tenggelam di bawah tetangga-tetangganya:
zona subduksi (\cref{ch:g11:subduction-zones}).

\textbf{11.} $d(1) = \qty{2850}{m}$, $d(4) = 2500 + 700 = \qty{3200}{m}$,
$d(16) = 2500 + 1400 = \qty{3900}{m}$, $d(64) = 2500 + 2800 = \qty{5300}{m}$.

\textbf{12.} $3900 - 2500 = \qty{1400}{m}$ selama 16 juta tahun pertama,
dan $5300 - 3900 = \qty{1400}{m}$ antara 16 dan 64 juta tahun: penenggelaman yang
sama, dalam waktu tiga kali lebih lama. Dasar itu tenggelam makin lama makin
lambat seiring umurnya.

\textbf{13.} $350\sqrt{t} = 1500$ memberikan $\sqrt{t} \approx 4.29$, jadi
$t \approx 18.4$ juta tahun.

\textbf{14.} $6400 - 3200\,e^{-64/62.8} \approx 6400 - 3200 \times 0.361
\approx \qty{5245}{m}$, sekitar \qty{55}{m} kurang dari hukum akar kuadrat: kedua
hukum itu sesuai sampai sekitar \qty{1}{\percent}.

\textbf{15.} Median yang terukur sekitar \qty{650}{m} lebih dangkal daripada
hukum itu. Hukum itu memerikan dasar yang dikoreksi untuk berat sedimennya;
kedalaman yang terukur mencakup sedimen, yang terkumpul selama puluhan juta tahun
di dasar seperti itu, serta gunung bawah laut dan dataran tinggi vulkanik, yang
semuanya membuat dasar yang terukur lebih dangkal.

\textbf{16.} $\alpha\,\Delta T = 3.2 \times 10^{-5} \times 700 = 0.0224$: bagian
mantelnya \qty{2.24}{\percent} lebih rapat,
$3380 \times 1.0224 \approx \qty{3456}{kg/m^3}$.

\textbf{17.} Bagian mantelnya setebal $100 - 7.1 = \qty{92.9}{km}$. Massa jenis
rata-rata: $(7.1 \times 2900 + 92.9 \times 3456)/100 \approx (20590 + 321060)/100
\approx \qty{3416}{kg/m^3}$.

\textbf{18.} $3416 > 3380$: \omterm{def:g11:lithosphere-asthenosphere:lithosphere}{litosfer} tua itu secara keseluruhan sekitar
\qty{36}{kg/m^3}, sekitar \qty{1}{\percent}, lebih rapat daripada \omterm{def:g11:lithosphere-asthenosphere:lithosphere}{astenosfer} di
bawahnya. Kerak yang ringan tidak lagi cukup untuk mengimbangi bagian mantelnya
yang dingin dan rapat.

\textbf{19.} Karena lebih rapat daripada \omterm{def:g11:lithosphere-asthenosphere:lithosphere}{astenosfer}, \omterm{def:g11:lithosphere-asthenosphere:lithosphere}{litosfer} itu dapat tenggelam
kembali ke dalam mantel di palung: subduksi, pokok bahasan bab berikutnya.

\textbf{20.} Pada umur 64 juta tahun potongan dasar itu terletak pada kedalaman
sekitar \qty{5300}{m}: sekitar \textbf{5.3~km} di bawah permukaan laut,
\qty{2.8}{km} lebih dalam daripada tempat kelahirannya.
\end{solution}
